\chapter{Khởi Động Sau Giờ Ra Chơi}
\chapterintro{Kéo lớp từ sân vào tiết}{Đây là thời điểm học sinh còn đầy chuyện dang dở, nên khởi động phải vừa mềm vừa dứt.}

\section{Một câu khóa từ sân vào lớp}
\begin{khoidongbox}{Hoạt động khởi động}
Nói một câu khóa rõ ràng: \textit{``Giờ chơi dừng ở cửa lớp, giờ học bắt đầu từ đây.''}
\end{khoidongbox}
\begin{visaohieuqua}
Học sinh cần một dấu mốc chuyển trạng thái. Câu khóa ngắn giúp cắt gọn mà không làm không khí nặng.
\end{visaohieuqua}
\begin{cachlam5phut}
Nói đúng một lần, sau đó cho lớp làm ngay một việc rất nhỏ như nhìn lên bảng hoặc trả lời một câu dễ.
\end{cachlam5phut}
\ngancach

\section{Hít thở hai nhịp rồi vào bài}
\begin{khoidongbox}{Hoạt động khởi động}
Mời cả lớp hít vào, thở ra hai nhịp thật ngắn trước khi mở bài.
\end{khoidongbox}
\begin{visaohieuqua}
Sau giờ ra chơi, lớp thường còn dư nhịp chạy và nói. Hai nhịp thở là cách kéo tốc độ xuống vừa đủ để học.
\end{visaohieuqua}
\begin{cachlam5phut}
Không biến thành bài tập thư giãn dài. Làm xong thì chốt ngay: \textit{``Rồi, bắt đầu.''}
\end{cachlam5phut}
\ngancach

\section{Một từ bỏ ngoài cửa}
\begin{khoidongbox}{Hoạt động khởi động}
Hỏi vui: \textit{``Nếu phải bỏ lại ngoài cửa một thứ để dễ học hơn, em bỏ gì?''}
\end{khoidongbox}
\begin{visaohieuqua}
Câu hỏi này làm học sinh tự nhận ra điều đang kéo mình khỏi tiết học. Nó vừa hài hước vừa có tác dụng thật.
\end{visaohieuqua}
\begin{cachlam5phut}
Lấy 3 câu trả lời nhanh như: nói chuyện dở, cơn buồn ngủ, chuyện ở sân. Sau đó kéo thẳng vào bài.
\end{cachlam5phut}
\ngancach

\section{Giữ lại điều gì từ giờ chơi?}
\begin{khoidongbox}{Hoạt động khởi động}
Đôi khi giờ chơi có năng lượng tốt. Hỏi: \textit{``Điều gì từ giờ chơi mình muốn mang vào tiết này?''}
\end{khoidongbox}
\begin{visaohieuqua}
Không phải lúc nào cũng phải cắt sạch năng lượng cũ. Có thể giữ lại sự vui vẻ hoặc tinh thần đồng đội để vào bài tốt hơn.
\end{visaohieuqua}
\begin{cachlam5phut}
Mời 2 em nói, rồi chốt: \textit{``Mình giữ tiếng cười, bỏ phần ồn.''}
\end{cachlam5phut}
\ngancach

\section{Mở bài bằng chuyện vừa cười}
\begin{khoidongbox}{Hoạt động khởi động}
Chọn một câu chuyện vui rất ngắn vừa xảy ra ở giờ chơi và nối nó sang chủ đề bài học.
\end{khoidongbox}
\begin{visaohieuqua}
Lớp vào bài nhanh hơn khi thấy tiết học không tách rời hoàn toàn khỏi đời sống ngay trước đó. Nhưng phần nối phải thật gọn.
\end{visaohieuqua}
\begin{cachlam5phut}
Kể trong 15 giây, không sa đà. Ngay sau đó nói rõ: \textit{``Từ chuyện này, mình vào bài hôm nay.''}
\end{cachlam5phut}